\documentclass[11pt,letterpaper]{article}
\usepackage[letterpaper,textwidth=437pt,textheight=634pt,hcentering,top=81pt]{geometry}
\usepackage[T1]{fontenc}
\usepackage{lmodern,amsmath,amssymb,amsthm}
\usepackage[hidelinks]{hyperref}
\hypersetup{pdftitle={A kernel for short solutions of the generalized Rubik's Cube},pdfauthor={Ruimin Yan}}
\newtheorem{theorem}{Theorem}
\newtheorem{lemma}[theorem]{Lemma}
\newtheorem{corollary}[theorem]{Corollary}
\title{\bfseries A kernel for short solutions of the generalized Rubik's Cube}
\author{Ruimin Yan}
\date{November 26, 2026\\[3pt]\small Working manuscript}
\begin{document}
\maketitle
\begin{abstract}
Given a six-colour cube of side $n$ and a budget of $k$ single-layer turns,
we give a polynomial-time reduction that either rejects the instance or produces
an equivalent cube of the same parity and side at most $4k+7$.
The output uses at most $6(4k+7)^2$ stickers. The same reduction works for
a prescribed solved colouring and for the goal that each face has one colour
and all six colours differ. A lifting argument uses two spare mirror depths
to verify every omitted sticker. The reduction, equivalence of the physical
instances, and a search bound have been formalized in Lean~4.
\end{abstract}

\section{The decision problem}
Each move rotates one layer orthogonal to a coordinate axis by
$90^\circ$, $-90^\circ$, or $180^\circ$. Each costs one move. The input lists
the colours on the six $n$ by $n$ faces in fixed coordinates, with $n\ge2$.
Demaine, Eisenstat and Rudoy proved NP-completeness of the unrestricted
decision problem \cite{der}. Their colour goal permits any six distinct
monochromatic faces. We study the dependence of exact search on the move
budget. Direct enumeration already decides the problem in at most
$\sum_{i=0}^k(9n)^i$ candidate words. Its dependence on $n$ motivates
reducing the set of candidate layers before searching.

\begin{theorem}\label{thm:main}
Given a legal coloured cube $X$ and an integer $k\ge0$, a polynomial-time
algorithm certifies that the distance to the prescribed solved colouring
exceeds $k$, or returns a legal cube $Y$ of side $2\le n'\le4k+7$
and of the same parity as $n$. The original instance satisfies $d(X)\le k$
if and only if $d(Y)\le k$. The same choice of retained layers preserves solvability within $k$ moves
for the goal of six distinctly coloured monochromatic faces.
\end{theorem}

For example, with $n=1000$ and $k=10$, every returned instance has even
side at most $46$. The decision concerns this particular budget. The
argument makes no assertion about equality of distances exceeding it.

\section{Mirror depths and changes of coordinates}
Number the layers on each axis from $0$ through $n-1$. The two outer
layers are distinguished, as is the central layer when $n$ is odd. Group
the remaining indices in pairs $\{i,n-1-i\}$, called ordinary depths.
There are $r=\lfloor(n-2)/2\rfloor$ such pairs. Use the same depth labels
on all three axes. A permutation of the $r$ labels acts on the layers of
each axis, preserving the side within each pair. A surface sticker has a
face and two tangential coordinates. Its support, the set of ordinary
depth labels among those coordinates, has size at most two.

A relabelling commutes with coordinate reversal $i\mapsto n-1-i$. A
quarter turn of a layer tests its layer index and exchanges its two tangential
coordinates with one reversal. Relabelling therefore carries that turn to the
turn on the renamed layer. It fixes the prescribed solved colouring, which
depends only on the face. These identities extend to turn words by composition.
A word of length at most $k$ names at most $k$ ordinary depths.

Call two depths equivalent for $X$ if their transposition, applied
simultaneously on all three axes, preserves every sticker colour of $X$.
This is an equivalence relation. Indeed the stabilizer of $X$ is a group,
and conjugating $(b\,c)$ by $(a\,b)$ gives $(a\,c)$ for distinct labels.
If a word of length at most $k$ generates $X$ from the prescribed target,
all unused depths belong to one class. Swapping any two of them fixes both
the word and the solved colouring. Inversion exchanges generating and
solving words and preserves length.

\section{Reduction and proof}
If $r\le2k+2$, retain all depths. Otherwise find an equivalence class $C$
whose complement has size at most $k$. If none exists, the unused-depth
observation certifies rejection. Keep all of the complement and exactly
$k+2$ members of $C$. This is possible since $r>2k+2$. Let $R$ be the
retained set. Restrict each face to the surviving rows and columns and
renumber the depths consecutively in their original order. Keep both outer
layers and the odd central layer.

\begin{lemma}\label{lem:lift}
The restricted colouring can be generated from the prescribed target in
at most $k$ turns if and only if the original colouring can be generated
from its prescribed target in at most $k$ turns.
\end{lemma}
\begin{proof}
Any word on the original cube projects by omitting turns on removed layers.
A retained sticker stays in retained coordinates under every retained turn.
A removed-layer turn fixes it. Thus projection preserves the colour read
at each retained sticker and cannot increase the length. Projection also
preserves legality for an input generated by an arbitrarily long word.

Conversely, embed a short word on the smaller cube by replacing each new
depth label with its original label. Call the colouring it produces $Z$.
The turn identities above give $Z=X$ on all retained stickers. If all
depths were retained we are done. Otherwise at least two of the $k+2$
retained members of $C$ are unused by the embedded word. For an omitted
sticker choose a distinct spare representative for each omitted depth in
its support, avoiding the retained depths already in that support.
There are enough spares because the support has size at most two.
Swap each omitted label with its chosen representative. Every such swap
fixes $X$ by the definition of $C$, and fixes $Z$ because its labels are
unused by the embedded word. Together the swaps carry the sticker to the
retained set. Equality at its image proves equality at the original
sticker. This proves $Z=X$ everywhere.
\end{proof}

\begin{proof}[Proof of Theorem~\ref{thm:main}]
There are at most $2k+2$ retained depths. With two outer layers and
possibly one central layer, the side is at most $2(2k+2)+3=4k+7$.
The parity is unchanged. Lemma~\ref{lem:lift} gives both directions.
For the monochromatic goal, a completed cube assigns the six colours
bijectively to the six faces. A bijection on colours preserves every
equality tested by a depth swap. All $6!$ target assignments consequently
use the same retained set. Applying the fixed-target equivalence to each
assignment proves the second assertion.
\end{proof}

\section{Search complexity and verification}
With random access to the six colour arrays, the full table testing
ordered depth pairs has $6r^2n^2$ colour comparisons. Selecting a class,
sorting representatives and materializing the smaller faces take polynomial
time in the input size. The smaller cube has $9n'$ moves. For a fixed target
the search considers at most $(9(4k+7)+1)^k$ candidate words, including
prefixes. Checking a word against the smaller colouring costs a polynomial
in $k+1$. This yields a fixed-parameter algorithm. Testing all permitted
face colours multiplies the search bound by at most $720$.

The physical decision theorem is in \texttt{CubeKernel.Monochrome}, the six-array encoding in
\texttt{CubeKernel.Packed}, and the search-tree and comparison counts in
\texttt{CubeKernel.Cost}. For a prescribed target, the implementation in
\texttt{CubeKernel.Solution} returns solving turns with the original layer
numbers. The download includes the Lean toolchain and dependency lock.
The complexity statement uses the array model above. The checked
counts bound comparisons in the complete table and words in the search tree.

\paragraph{Status and assistance.}
This working manuscript reports a machine-checked result. The literature
comparison does not establish a claim of priority. OpenAI Codex assisted in
developing the argument, implementing the Lean proofs and drafting the text.
\begin{thebibliography}{9}
\bibitem{der} E.~D. Demaine, S.~Eisenstat and M.~Rudoy.
\href{https://erikdemaine.org/papers/Rubik_STACS2018/paper.pdf}{Solving the Rubik's Cube optimally is NP-complete}.
35th Symposium on Theoretical Aspects of Computer Science, 2018, article~25.
\end{thebibliography}
\end{document}
